%-----------------------------------------------------------------------------------------------------------------------------------------------%
%	The MIT License (MIT)
%
%	Copyright (c) 2015 Jan Küster
%
%	Permission is hereby granted, free of charge, to any person obtaining a copy
%	of this software and associated documentation files (the "Software"), to deal
%	in the Software without restriction, including without limitation the rights
%	to use, copy, modify, merge, publish, distribute, sublicense, and/or sell
%	copies of the Software, and to permit persons to whom the Software is
%	furnished to do so, subject to the following conditions:
%	
%	THE SOFTWARE IS PROVIDED "AS IS", WITHOUT WARRANTY OF ANY KIND, EXPRESS OR
%	IMPLIED, INCLUDING BUT NOT LIMITED TO THE WARRANTIES OF MERCHANTABILITY,
%	FITNESS FOR A PARTICULAR PURPOSE AND NONINFRINGEMENT. IN NO EVENT SHALL THE
%	AUTHORS OR COPYRIGHT HOLDERS BE LIABLE FOR ANY CLAIM, DAMAGES OR OTHER
%	LIABILITY, WHETHER IN AN ACTION OF CONTRACT, TORT OR OTHERWISE, ARISING FROM,
%	OUT OF OR IN CONNECTION WITH THE SOFTWARE OR THE USE OR OTHER DEALINGS IN
%	THE SOFTWARE.
%	
%
%-----------------------------------------------------------------------------------------------------------------------------------------------%


%============================================================================%
%
%	DOCUMENT DEFINITION
%
%============================================================================%

%we use article class because we want to fully customize the page and dont use a cv template
\documentclass[10pt,A4]{article}	


%----------------------------------------------------------------------------------------
%	ENCODING
%----------------------------------------------------------------------------------------

%we use utf8 since we want to build from any machine
\usepackage[utf8]{inputenc}		

\usepackage{CJKutf8}

%----------------------------------------------------------------------------------------
%	LOGIC
%----------------------------------------------------------------------------------------

% provides \isempty test
\usepackage{xifthen}

\usepackage{xifthen}

\usepackage{color}

%accent color
\definecolor{sectcol}{RGB}{0,150,255}

%dark background color
\definecolor{bgcol}{RGB}{110,110,110}

%light background / accent color
\definecolor{softcol}{RGB}{225,225,225}


\definecolor{linkcol}{RGB}{0,150,255}


\usepackage{hyperref}
\hypersetup{backref,
  pdftitle=Geza Kovacs CV,
  pdfauthor=Geza Kovacs,
  pdfkeywords=Geza Kovacs CV,
  pdfsubject=Geza Kovacs CV,
  colorlinks=true,
  urlcolor=linkcol
}

% \makeatletter
% \def\url@leostyle{%
%   \@ifundefined{selectfont}{\def\UrlFont{\sf}}{\def\UrlFont{\bf\ttfamily}}}
% \makeatother
% \urlstyle{leo}

\urlstyle{sf}
\makeatletter
% Inspired by http://anti.teamidiot.de/nei/2009/09/latex_url_slash_spacingkerning/
% but slightly less kern and shorter underscore
\let\UrlSpecialsOld\UrlSpecials
\def\UrlSpecials{\UrlSpecialsOld\do\/{\Url@slash}\do\_{\Url@underscore}}%
\def\Url@slash{\@ifnextchar/{\kern-.11em\mathchar47\kern-.2em}%
    {\kern-.0em\mathchar47\kern-.08em\penalty\UrlBigBreakPenalty}}
\def\Url@underscore{\nfss@text{\leavevmode \kern.06em\vbox{\hrule\@width.3em}}}
\makeatother

\newcommand{\bhref}[2]{\href{#1}{\bfseries #2}}

%\newcommand{\burl}[1]{\href{#1}{\bfseries #1}}
\newcommand{\burl}[1]{\begin{bfseries}\url{#1}\end{bfseries}}

%----------------------------------------------------------------------------------------
%	FONT
%----------------------------------------------------------------------------------------

% some tex-live fonts - choose your own

%\usepackage[defaultsans]{droidsans}
%\usepackage[default]{comfortaa}
%\usepackage{cmbright}
\usepackage[default]{raleway}
%\usepackage{fetamont}
%\usepackage[default]{gillius}
%\usepackage[light,math]{iwona}
%\usepackage[thin]{roboto} 

% set font default
\renewcommand*\familydefault{\sfdefault} 	
\usepackage[T1]{fontenc}

% more font size definitions
\usepackage{moresize}		


%----------------------------------------------------------------------------------------
%	PAGE LAYOUT  DEFINITIONS
%----------------------------------------------------------------------------------------

%debug page outer frames
%\usepackage{showframe}			


%define page styles using geometry
\usepackage[a4paper]{geometry}		

% for example, change the margins to 2 inches all round
%\geometry{top=1.25cm, bottom=-.6cm, left=1.5cm, right=1.5cm} 	
\geometry{top=1.25cm, bottom=1.25cm, left=1.5cm, right=1.5cm} 	

%use customized header
\usepackage{fancyhdr}				
\pagestyle{fancy}

%less space between header and content
\setlength{\headheight}{-5pt}		


%customize entries left, center and right
%\lhead{}
%\chead{}
%\rhead{}


%indentation is zero
\setlength{\parindent}{0mm}

%----------------------------------------------------------------------------------------
%	TABLE /ARRAY DEFINITIONS
%---------------------------------------------------------------------------------------- 

%for layouting tables
\usepackage{multicol}			
\usepackage{multirow}

%extended aligning of tabular cells
\usepackage{array}

\newcolumntype{x}[1]{%
>{\raggedleft\hspace{0pt}}p{#1}}%


%----------------------------------------------------------------------------------------
%	GRAPHICS DEFINITIONS
%---------------------------------------------------------------------------------------- 

%for header image
\usepackage{graphicx}

%for floating figures
\usepackage{wrapfig}
\usepackage{float}
%\floatstyle{boxed} 
%\restylefloat{figure}

%for drawing graphics		
\usepackage{tikz}				
\usetikzlibrary{shapes, backgrounds,mindmap, trees}


%----------------------------------------------------------------------------------------
%	Color DEFINITIONS
%---------------------------------------------------------------------------------------- 

%============================================================================%
%
%
%	DEFINITIONS
%
%
%============================================================================%

%----------------------------------------------------------------------------------------
% 	HEADER
%----------------------------------------------------------------------------------------

% remove top header line
\renewcommand{\headrulewidth}{0pt} 

%remove botttom header line
\renewcommand{\footrulewidth}{0pt}	  	

%remove pagenum
\renewcommand{\thepage}{}	

%remove section num		
\renewcommand{\thesection}{}			

%----------------------------------------------------------------------------------------
% 	ARROW GRAPHICS in Tikz
%----------------------------------------------------------------------------------------

% a six pointed arrow poiting to the left
\newcommand{\tzlarrow}{(0,0) -- (0.2,0) -- (0.3,0.2) -- (0.2,0.4) -- (0,0.4) -- (0.1,0.2) -- cycle;}	

% include the left arrow into a tikz picture
% param1: fill color
%
\newcommand{\larrow}[1]
{\begin{tikzpicture}[scale=0.58]
	 \filldraw[fill=#1!100,draw=#1!100!black]  \tzlarrow
 \end{tikzpicture}
}

% a six pointed arrow poiting to the right
\newcommand{\tzrarrow}{ (0,0.2) -- (0.1,0) -- (0.3,0) -- (0.2,0.2) -- (0.3,0.4) -- (0.1,0.4) -- cycle;}

% include the right arrow into a tikz picture
% param1: fill color
%
\newcommand{\rarrow}
{\begin{tikzpicture}[scale=0.7]
	\filldraw[fill=sectcol!100,draw=sectcol!100!black] \tzrarrow
 \end{tikzpicture}
}



%----------------------------------------------------------------------------------------
%	custom sections
%----------------------------------------------------------------------------------------

% create a coloured box with arrow and title as cv section headline
% param 1: section title
%
\newcommand{\cvsection}[1]
{
	\begin{center}
		\large\textcolor{sectcol}{\textbf{#1}}
	\end{center}
}

%create a coloured arrow with title as cv meta section section
% param 1: meta section title
%
\newcommand{\metasection}[2]
{
%\begin{tabular*}{1\textwidth}{r r}
\footnotesize{#2} \hspace*{\fill} \footnotesize{#1}\\[1pt]
%\end{tabular*}
}

%----------------------------------------------------------------------------------------
%	 CV EVENT
%----------------------------------------------------------------------------------------

% creates a stretched box as cv entry headline followed by two paragraphs about 
% the work you did
% param 1:	event time i.e. 2014 or 2011-2014 etc.
% param 2:	event name (what did you do?)
% param 3:	institution (where did you work / study)
% param 4:	what was your position
% param 5:	some words about your contributions
%
% \newcommand{\cvevent}[5]
% {

% \begin{tabular*}{1\textwidth}{p{13.6cm}  x{3.9cm}}
% 	\textbf{#2} - \textcolor{bgcol}{#3} &   \vspace{2.5pt}\textcolor{sectcol}{#1}
% \end{tabular*}

% \vspace{-8pt}
% \textcolor{softcol}{\hrule}
% \vspace{6pt}

%   $\cdot$ #4\\[3pt]
%   $\cdot$ #5\\[6pt]

% }

\newcommand{\cvevent}[5]
{

	\textbf{#2} - \textcolor{bgcol}{#3} \hfill \vspace{2.5pt}\textcolor{sectcol}{#1}

\vspace{-8pt}
\textcolor{softcol}{\hrule}
\vspace{6pt}

  #4\\[3pt]
  #5\\[6pt]

}

\newcommand{\cvevtzero}[4]
{

	\textbf{#2} - \textcolor{bgcol}{#3} \hfill \vspace{2.5pt}\textcolor{sectcol}{#1}

\vspace{-8pt}
\textcolor{softcol}{\hrule}
\vspace{6pt}

}

% \newcommand{\cvevtzero}[3]
% {

% \begin{tabular*}{1\textwidth}{p{13.6cm}  x{3.9cm}}
% 	\textbf{#2} - \textcolor{bgcol}{#3} &   \vspace{2.5pt}\textcolor{sectcol}{#1}
% \end{tabular*}

% \vspace{-8pt}
% \textcolor{softcol}{\hrule}
% \vspace{6pt}
% }

% \newcommand{\cvevtone}[4]
% {

% \begin{tabular*}{1\textwidth}{p{13.6cm}  x{3.9cm}}
% 	\textbf{#2} - \textcolor{bgcol}{#3} &   \vspace{2.5pt}\textcolor{sectcol}{#1}
% \end{tabular*}

% \vspace{-8pt}
% \textcolor{softcol}{\hrule}
% \vspace{6pt}

%   $\cdot$ #4\\[3pt]

% }

\newcommand{\cvevtone}[4]
{

	\textbf{#2} - \textcolor{bgcol}{#3} \hfill \vspace{2.5pt}\textcolor{sectcol}{#1}

\vspace{-8pt}
\textcolor{softcol}{\hrule}
\vspace{6pt}

  #4\\[3pt]

}


% creates a stretched box as 
\newcommand{\cveventmeta}[2]
{
	\mbox{\mystrut \hspace{87pt}\textit{#1}}\\
	#2
}

%----------------------------------------------------------------------------------------
% CUSTOM STRUT FOR EMPTY BOXES
%----------------------------------------- -----------------------------------------------
\newcommand{\mystrut}{\rule[-.3\baselineskip]{0pt}{\baselineskip}}

%----------------------------------------------------------------------------------------
% CUSTOM LOREM IPSUM
%----------------------------------------------------------------------------------------
\newcommand{\lorem}
{Lorem ipsum dolor sit amet, consectetur adipiscing elit. Donec a diam lectus.}



%============================================================================%
%
%
%
%	DOCUMENT CONTENT
%
%
%
%============================================================================%
\begin{document}


%use our custom fancy header definitions
\pagestyle{fancy}	





%----------------------------------------------------------------------------------------
%	HEADER IMAGE
%----------------------------------------------------------------------------------------

%\begin{figure}[H]
%\begin{flushright}
%	\includegraphics[trim= 320 130 460 210,clip,width=0.2\linewidth]{myfoto.jpg}	%trimming relative to image size!
%\end{flushright}
%\end{figure}


%---------------------------------------------------------------------------------------
%	TITLE HEADLINE
%----------------------------------------------------------------------------------------
\vspace{-8pt}
\begin{center}
	\HUGE \textsc{Geza Kovacs}\\[2pt]
	\small \hypersetup{urlcolor=black}\href{https://www.gkovacs.com}{www.gkovacs.com}\hypersetup{urlcolor=linkcol} \hfill \hspace{5mm} Human-Computer Interaction Researcher \hfill \hypersetup{urlcolor=black}\href{mailto:geza@cs.stanford.edu}{geza@cs.stanford.edu}\hypersetup{urlcolor=linkcol}
	%\small www.gkovacs.com
\end{center}



%\vspace{6pt}


%---------------------------------------------------------------------------------------
%	META SECTION
%----------------------------------------------------------------------------------------
%\metasection{Bremen, Germany}{\textbf{Status:} M.Sc. Digital Media, Scientific Employee}
%\metasection{www.gkovacs.com}{\textbf{Fields:} Software Engineering, Consulting, Scrum, Usability} 
%\metasection{info@jankuester.com}{\textbf{Techs:} Meteor, Javascript, Mongo,. Bootstrap, Bash, Git, Webstorm}
%\metasection{+49 176 *** *** **}{\textbf{Activities:} Global Game Jam, Open Source, Fitness}
%\vspace{-2pt}
%\textcolor{softcol}{\hrule}
%\vspace{6pt}

\normalsize

%---------------------------------------------------------------------------------------
%	SUMMARAY (optional)
%----------------------------------------------------------------------------------------
%\vspace{-6pt}
%\cvsection{Research Interests}
%Currently working on human-in-the-loop machine translation systems, and computer-assisted translation. PhD work was on behavior change.
%Digital media graduate with project experience in the field of technology based assessment. Currently working as IT Consultant at We4IT, Bremen in the field of IBM Notes Domino and XPages applications. Master studies focused on teams from different disciplines and cultural backgrounds on solutions for complex problems.\\



%---------------------------------------------------------------------------------------
%	EDUCATION SECTION
%--------------------------------------------------------------------------------------
\cvsection{Education}

\cvevtone{2019}{Ph.D, Computer Science}{Stanford University}{Thesis: HabitLab: In-the-wild Behavior Experimentation at Scale. Advisor: Michael Bernstein. \hfill GPA: 4.0/4.0}%{Developed and evaluated an algorithm for semi automated scoring of spreadsheet data}

%\textcolor{softcol}{\hrule}

%
\cvevtone{2013}{B.S. and M.Eng, Computer Science}{Massachusetts Institute of Technology}{Thesis: Multimedia for Language Learning. Advisor: Robert Miller. \hfill GPA: 5.0/5.0}

%\textcolor{softcol}{\hrule}

%
%\cvevent{2012 - 2015}{Master Studies Digital Media}{University of Bremen}{Inter-cultural classes in English, covering special topics in computer science and design}{Professionalized in research methods, software development and e-assessment}

%\textcolor{softcol}{\hrule}

%
%\cvevent{2009 - 2010}{Semester Abroad}{University of Melbourne}{Mastered six months of study and trans-cultural experience in Melbourne, Australia}{Finished machine programming, information visualization, professional essay writing}


%============================================================================%
%
%	CV SECTIONS AND EVENTS (MAIN CONTENT)
%
%============================================================================%

%---------------------------------------------------------------------------------------
%	EXPERIENCE
%----------------------------------------------------------------------------------------
\cvsection{Experience}

\cvevtzero{2019 - present}{Senior Research Scientist}{Lilt, San Francisco. Manager: John DeNero}\\
Lilt is a human-in-the-loop machine translation startup. I lead the HCI research group, which consists of 3 FTEs.\\
Improved interactive MT system speed by shifting computation client-side and reusing computation via heuristics.\\
Ran user studies to evaluate where translators spend their time during translation, and identify features of translator behaviors that predict better performance. Ran A/B tests to evaluate the return-on-investment of translating websites, and developed a data-driven system to recommend which pages to translate into which languages.\\

\cvevtzero{2013 - 2019}{Ph.D Student}{Stanford University. Advisor: Michael Bernstein}\\
My thesis project was HabitLab \url{http://habitlab.stanford.edu/} an app for Chrome + Android with \textbf{12,000+ daily active users} which helps users reduce time online. Published first-author papers about HabitLab at CHI 2021, CHI 2019, and CSCW 2018. I also published first-author papers at Learning at Scale 2016 (Effects of In-Video Quizzes on MOOC Lecture Viewing) and CHI 2014 (Smart Subtitles for Foreign Language Learning).\\ % See page 2 for other publications.

%
\cvevtone{Summer 2015}{Research Intern}{Microsoft Research, Redmond. Manager: Merrie Morris}{Designed and built an educational social feed app usable by pre-literate children. Published at CSCW 2017.}

\cvevtone{Summer 2014}{Research Intern}{Microsoft Research, Beijing. Manager: Darren Edge}{Built QuizCram, a quiz-driven MOOC lecture viewer that improves learning outcomes. Published at CHI 2015 SRC.}

\cvevtone{Summer 2013}{Research Intern}{Google Research, Mountain View. Manager: Shumin Zhai}{Developed a machine learning system for detecting taps on the phone bezel, for use in Android input methods.}

\cvevtone{Summer 2012}{Software Engineering Intern}{Google, Mountain View}{Developed an NLP system to detect vocabulary and generate glossaries from book text (used MapReduce).}

\cvevtone{Summer 2011}{Software Engineering Intern}{Google, Mountain View}{Developed a machine learning system to predict the quality of user reviews, now deployed on Google Play.}

\cvevtone{Summer 2010}{Software Development Engineering Intern}{Microsoft, Redmond}{Built the intellsense code completion system for a scientific computing language, and contributed to its compiler.}

%--------------------------------------------------------------------------------------------------
%	ARTIFICIAL FOOTER (fancy footer cannot exceed linewidth) 
%--------------------------------------------------------------------------------------------------

\cvsection{Open-source Projects}

% \url{http://unetbootin.github.io/}
\textbf{UNetbootin (LiveUSB Creator)} \hfill \url{https://en.wikipedia.org/wiki/UNetbootin}\\ %\hfill \textsl{January 2007 -- now}\\
\emph{40 million downloads}. UNetbootin creates bootable USB flash drives for various (50+) Linux distributions.\\ %This work has been accepted into the official package repositories for Debian, Ubuntu, Fedora, openSUSE, Gentoo, and other major distributions. \\
% \hspace{3mm}  

\textbf{Ubuntu Installer for Windows (Wubi)} \hfill \url{https://en.wikipedia.org/wiki/Wubi_(software)}\\ %\hfill \textsl{November 2006 -- August 2007}\\
\emph{Now part of Ubuntu}. Built the first versions of Wubi, which allows Ubuntu to be installed from Windows.\\ % This work is now part of Ubuntu. \\

\textbf{HabitLab (In-the-wild Behavior Change Research Platform)} \hfill \url{https://habitlab.stanford.edu}\\ %\hfill \textsl{November 2006 -- August 2007}\\
\emph{12,000+ daily active users}. I built HabitLab over my Ph.D, and it is still used for research at Stanford Medical School.\\ % This work is now part of Ubuntu. \\


\cvsection{Conference and Journal Papers}

\textbf{Geza Kovacs}, Zhengxuan Wu, Michael Bernstein. ``Not Now, Ask Later: Users Weaken Their Behavior Change Regimen Over Time, But Expect To Re-Strengthen It Imminently'' ACM annual conference on Human Factors in Computing Systems (CHI) 2021 (to appear).\\ %Acceptance rate: 23.8\%

Samuel Läubli, Patrick Simianer, Joern Wuebker, \textbf{Geza Kovacs}, Rico Sennrich, Spence Green. ``The Impact of Text Presentation on Translator Performance.'' Target: International Journal of Translation Studies, 2021 (to appear).\\

\textbf{Geza Kovacs}, Drew Mylander Gregory, Zilin Ma, Zhengxuan Wu, Golrokh Emami, Jacob Ray, Michael Bernstein. ``Conservation of Procrastination: Do Productivity Interventions Save Time Or Just Redistribute It?'' ACM annual conference on Human Factors in Computing Systems (CHI) 2019.\\ %Acceptance rate: 23.8\%

\textbf{Geza Kovacs}, Zhengxuan Wu, Michael Bernstein. ``Rotating Online Behavior Change Interventions Increases Effectiveness But Also Increases Attrition.'' ACM annual conference on Computer-Supported Cooperative Work and Social Computing (CSCW) 2018.\\ %Acceptance rate: 26\%

Rajan Vaish, Neil Gaikwad, \textbf{Geza Kovacs}, Andreas Veit, Ranjay Krishna, Imanol Arrieta Ibarra, Camelia Simoiu, Michael Wilber, Serge Belongie, Sharad Goel, James Davis, Michael Bernstein. ``Crowd Research: Open and Scalable University Laboratories.'' ACM Symposium on User Interface Software and Technology (UIST) 2017.\\ %Acceptance rate: 22\%

Kiley Sobel, \textbf{Geza Kovacs}, Galen McQuillen, Andrew Cross, Nirupama Chandrasekaran, Nathalie Riche, Ed Cutrell, Meredith Morris. ``EduFeed: A Social Feed to Engage Preliterate Children in Educational Activities.'' ACM annual conference on Computer-Supported Cooperative Work and Social Computing (CSCW) 2017.\\ %Acceptance rate: 35\%

\textbf{Geza Kovacs}. ``Effects of In-Video Quizzes on MOOC Lecture Viewing.'' ACM annual conference on Learning at Scale (L@S) 2016.\\ %Acceptance rate: 22\%

\textbf{Geza Kovacs} and Robert C. Miller. ``Smart Subtitles for Vocabulary Learning.'' ACM annual conference on Human Factors in Computing Systems (CHI) 2014.\\ %Acceptance rate: 23\%

\cvsection{Peer-Reviewed Extended Abstracts}

Stanford Crowd Research, \textbf{Geza Kovacs}, Rajan Vaish, Michael Bernstein. ``Daemo: A Self-Governed Crowdsourcing Marketplace''. ACM Symposium on User Interface Software and Technology (UIST) 2015, Poster.\\

\textbf{Geza Kovacs}. ``FeedLearn: Using Facebook Feeds for Microlearning.'' ACM annual conference on Human Factors in Computing Systems (CHI) 2015, Extended Abstracts.\\ % Acceptance rate: 45\%

\textbf{Geza Kovacs}. ``QuizCram: A Question-Driven Video Studying Interface.'' ACM annual conference on Human Factors in Computing Systems (CHI) 2015, Extended Abstracts.\\

Joseph Jay Williams, \textbf{Geza Kovacs}, Caren Walker, Samuel G Maldonado, Tania Lombrozo. ``Learning Online via Prompts to Explain.'' ACM annual conference on Human Factors in Computing Systems (CHI) 2014, Extended Abstracts.\\

\textbf{Geza Kovacs} and Robert C. Miller. ``Foreign Manga Reader: Learn Grammar and Pronunciation while Reading Comics.'' ACM Symposium on User Interface Software and Technology (UIST) 2013, Demo.\\

\textbf{Geza Kovacs}. ``Smart Subtitles for Language Learning.'' ACM annual conference on Human Factors in Computing Systems (CHI) 2013, Extended Abstracts.\\ %\\
%\url{http://groups.csail.mit.edu/uid/other-pubs/chi2013-smartsubs.pdf}

\textbf{Geza Kovacs}. ``ScreenMatch: providing context to software translators by displaying screenshots.'' ACM annual conference on Human Factors in Computing Systems (CHI) 2012, Extended Abstracts.\\ %\\
%\url{http://groups.csail.mit.edu/uid/other-pubs/chi2012-screenshots-for-translation-context.pdf}

%\vspace{-5mm}

\cvsection{Invited Talks}

\textbf{Geza Kovacs}. ``Predictive Translation Memory in the Wild: A Study of Interactive Machine Translation Use on Lilt.'' Association for Machine Translation in the Americas (AMTA) Workshop on the Impact of Machine Translation (iMpacT 2020).\\

\pagebreak

\cvsection{Awards and Honors}

Stanford Human-Centered AI Grant (for my research project HabitLab) \hfill \textcolor{sectcol}{2018}\\ %-2016\\
National Defense Science and Engineering Graduate Fellowship \hfill \textcolor{sectcol}{2013}\\ %-2016\\
National Science Foundation Graduate Research Fellowship \hfill \textcolor{sectcol}{2013}\\ %(declined in favor of NDSEG), 2013\\
% Finalist and Honorable Mention, MIT Web Programming Competition (6.470), 2013\\ % (Project: PsetParty)
%1$^{\textrm{st}}$ place, Most Useful, ACM UIST (User Interface Software and Technology) Student Innovation Contest \hfill \textcolor{sectcol}{2012}\\
1$^{\textrm{st}}$ place, ACM UIST (User Interface Software and Technology) Student Innovation Contest \hfill \textcolor{sectcol}{2012}\\
1$^{\textrm{st}}$ place, ACM CHI (Human Factors in Computing Systems) Student Research Competition \hfill \textcolor{sectcol}{2012}\\
%1$^{\textrm{st}}$ place award for my project, InstantKaraoke, at Boston Music Hack Day \hfill \textcolor{sectcol}{2012}\\ % (Project: InstantKaraoke)
Member of Phi Beta Kappa (Liberal Arts), Tau Beta Pi (Engineering), Eta Kappa Nu (EECS) honor societies \hfill \textcolor{sectcol}{2012}\\
1$^{\textrm{st}}$ place, MIT Autonomous Robotics Competition (MASLAB) \hfill \textcolor{sectcol}{2010}\\

\cvsection{Academic Service}

Reviewer, ACM Conference on Human Factors in Computing Systems (CHI) \hfill \textcolor{sectcol}{2015, 2018, 2019, 2021}

Reviewer, ACM Symposium on User Interface Software and Technology (UIST) \hfill \textcolor{sectcol}{2017, 2018}

Reviewer, ACM Conference on Computer-Supported Cooperative Work and Social Computing (CSCW) \hfill \textcolor{sectcol}{2021}

Reviewer, ACM Interactive, Mobile, Wearable and Ubiquitous Technologies (IMWUT) \hfill \textcolor{sectcol}{2019} \\


% Reviewer for \textbf{CHI} (2015, 2018, 2019, 2021), \textbf{UIST} (2017, 2018), \textbf{CSCW} (2021), \textbf{IMWUT} (2019).\\

\cvsection{Teaching Experience}

\textbf{Understanding Users (CS 377U)} -- Teaching Assistant, at Stanford \hfill \textcolor{sectcol}{Spring 2019} \\

\textbf{Human Computer Interaction Research (CS 376)} -- Teaching Assistant, at Stanford \hfill \textcolor{sectcol}{Fall 2018} \\
%Helped write and grade assignments, lead discussions, and manage the course infrastructure.

\textbf{Natural Language Processing (6.863)} -- Teaching Assistant, at MIT \hfill \textcolor{sectcol}{Fall 2012} \\
%Helped write and grade assignments, and managed the course infrastructure. I developed new tools to make the assignment grading process faster, semi-automatic, and paper-free.

\textbf{Introduction to C++ IAP (6.096)} -- Instructor, at MIT \hfill \textcolor{sectcol}{January 2011}

%Gave lectures, helped write and grade assignments, and helped students in lab for a student-run, for-credit introductory C++ course. The teaching materials I produced have been made available on OpenCourseWare:
My lectures and teaching materials for this course are available on MIT OpenCourseWare:

\begin{footnotesize}
\url{http://ocw.mit.edu/courses/electrical-engineering-and-computer-science/6-096-introduction-to-c-january-iap-2011}
\end{footnotesize}\\

\textbf{MASLAB Mobile Autonomous Systems Lab (6.186)} -- Software Director, at MIT \hfill \textcolor{sectcol}{January 2011}

Gave lectures on computer vision and managed the software stack for MIT's autonomous robotics competition.\\
% As the software director for this student-led autonomous robotics competition, I gave lectures on computer vision and control algorithms, managed the software stack for the competition, and helped students in lab.

\cvsection{Relevant Coursework}

Deep Learning (CS 230), Natural Language Processing (6.864, 6.863), Artificial Intelligence (6.034), Data Mining (CS 224w), Statistical Models (6.804), Statistics (18.440), Linear Algebra (18.700), Security (6.857), Bioinformatics (6.047), HCI (6.803), Algorithms (6.006, 6.046), Linguistics (24.900), Compilers (CS 143)\\

\cvsection{Selected Press}

\textcolor{sectcol}{\textbf{HabitLab}}\\

\textbf{WIRED} The HabitLab Browser Extension Curbs Your Time Wasted on the Web, January 2019.\\
\url{https://www.wired.com/story/habitlab-browser-extension/}\\

\textbf{Lifehacker} Prevent Procrastination With This Chrome Extension, February 2019.\\
\url{https://lifehacker.com/prevent-procrastination-with-this-chrome-extension-1832723418}\\

\textbf{The New York Times} Finding It Hard to Focus? Maybe It’s Not Your Fault, August 2018.\\
\url{https://www.nytimes.com/2018/08/14/style/how-can-i-focus-better.html}\\

\textbf{Lifehacker} Be More Mindful of the Time You Waste Online With HabitLab, August 2018.\\
\url{https://lifehacker.com/be-more-mindful-of-the-time-you-waste-online-with-habit-1828118354}\\

\textbf{WIRED} The Chrome Extensions We Can't Live Without, February 2018.\\
\url{https://www.wired.com/story/best-chrome-extensions/}\\

\textbf{How-To Geek} HabitLab Subtly Helps You Change Bad Online Habits, August 2018.\\
\url{https://www.howtogeek.com/fyi/free-download-habitlab-subtly-helps-you-change-bad-online-habits/}\\

\pagebreak

\textbf{The Stanford Daily} HabitLab browser extension aims to help users regain control of their online browsing behavior, March 2019.\\
\begin{footnotesize}
\url{https://www.stanforddaily.com/2019/03/13/habitlab-browser-extension-aims-to-help-users-regain-control-of-their-online-browsing-behavior/}
\end{footnotesize}\\

\vspace{-4mm}

\textbf{Entrepreneur} Use These Strategies to Maximize Productivity -- Without Inventing an Extra Weekday, May 2018.\\
\url{https://www.entrepreneur.com/article/312764}\\

%\textbf{PCMag} The Best Free Google Chrome Extensions, July 2020.\\
%\url{https://www.pcmag.com/news/the-100-best-free-google-chrome-extensions}\\

%\vspace{-4mm}

\textbf{Tencent News} (Chinese) \begin{CJK*}{UTF8}{gbsn}这款斯坦福大学的工具，让你远离加班，提升200％效率\end{CJK*}, March 2019.\\
\url{https://new.qq.com/omn/20190311/20190311A0BBRU.html}\\

%\pagebreak

\textcolor{sectcol}{\textbf{Crowd Research / Daemo}}\\

\textbf{Stanford University News} A Stanford-led platform for crowdsourced research gives experience to global participants, October 2017.\\
\url{https://news.stanford.edu/2017/10/23/crowdsourced-research-gives-experience-global-participants/}\\

\textbf{WIRED} Amazon's Turker Crowd Has Had Enough, August 2017.\\
\url{https://www.wired.com/story/amazons-turker-crowd-has-had-enough/}\\

\textcolor{sectcol}{\textbf{UNetbootin}}\\

\textbf{Forbes} How To Try Linux Without Making Any Changes To Your PC, September 2018.\\
\begin{small}
\url{https://www.forbes.com/sites/jasonevangelho/2018/09/18/how-to-safely-try-linux-on-your-mac-or-windows-pc/}
\end{small}\\

\textbf{PCWorld} Create a Bootable Linux Flash Drive in Three Easy Steps, February 2012.\\
\url{https://www.pcworld.com/article/249870/create_a_bootable_linux_flash_drive_in_three_easy_steps.html}\\

\textbf{Lifehacker} The Complete Guide to Saving Your Windows System with a Thumb Drive, March 2010.\\
\url{https://lifehacker.com/the-complete-guide-to-saving-your-windows-system-with-a-5504531}\\

%Installing Ubuntu on an old netbook with hair tearing and profanity. Network World, August 2014.\\
%\begin{footnotesize}
%\url{https://www.networkworld.com/article/2465150/installing-ubuntu-on-an-old-netbook-with-hair-tearing-and-profanity.html}
%\end{footnotesize}\\

%Linux on the NUC: Using Ubuntu, Mint, Fedora, and the SteamOS beta. Ars Technica, February 2014.\\
%\begin{small}
%\url{https://arstechnica.com/gadgets/2014/02/linux-on-the-nuc-using-ubuntu-mint-fedora-and-the-steamos-beta/}
%\end{small}\\

\textcolor{sectcol}{\textbf{Wubi}}\\

\textbf{Ars Technica} Wubi arrives: a look at Ubuntu 8.04 alpha 5, February 2008.\\
\url{https://arstechnica.com/information-technology/2008/02/wubi-arrives-a-look-at-ubuntu-8-04-alpha-5/}\\

\textbf{Lifehacker} Install Ubuntu on a Windows Netbook, No Partitioning Needed, May 2008.\\
\url{https://lifehacker.com/install-ubuntu-on-a-windows-netbook-no-partitioning-ne-5542387}\\

\textbf{PCWorld} The Ubuntu guide for displaced Windows users, March 2013.\\
\url{https://www.pcworld.com/article/2030132/the-ubuntu-guide-for-displaced-windows-users.html}\\

\cvsection{Management and Mentorship}

Sai Gouravajhala, Research Scientist at Lilt. \hfill \textcolor{sectcol}{August 2020 -- present}

Hannah Yan, Data Scientist at Lilt. \hfill \textcolor{sectcol}{September 2020 -- present}

Jessy Lin, Research Engineer at Lilt. Now a Ph.D student at UC Berkeley. \hfill \textcolor{sectcol}{August 2019 -- August 2020}

Ming-Chang Chiu, Data Science Intern at Lilt. Now a Ph.D student at USC. \hfill \textcolor{sectcol}{Summer 2020}

Zhengxuan Wu, MS student researcher at Stanford. \hfill \textcolor{sectcol}{Fall 2018 -- Spring 2019}

Zilin Ma, undergraduate researcher at Stanford. Now a Ph.D student at Harvard. \hfill \textcolor{sectcol}{Summer 2018}

Drew Mylander Gregory, undergraduate researcher at Stanford. \hfill \textcolor{sectcol}{Summer 2018}

Golrokh Emami, undergraduate researcher at Stanford. \hfill \textcolor{sectcol}{Summer 2018}

Jacob Ray, undergraduate researcher at Stanford. \hfill \textcolor{sectcol}{Summer 2018}\\

%\null
%\vspace*{\fill}
%\hspace{-0.25\linewidth}\colorbox{white}{\makebox[1.5\linewidth][c]{\mystrut  \textnormal{\textcolor{sectcol}{www.jankuester.com} $\cdot$ \textcolor{sectcol}{github.com/jankapunkt}}}}




%============================================================================%
%
%
%
%	DOCUMENT END
%
%
%
%============================================================================%
\end{document}
